\scanpage{131}

\mspage{468}

(\textsuperscript{107}) (4 octobre) J'ai eu occasion déjà de faire mention d'un aspect important de ces premières cinq années de ma vie, comme d'un"privilège" de grand prix (*): une identification profonde et sans problèmes avec mon père, laquelle n'a jamais été touchée par la peur ou par l'envie. Je me suis rendu compte de cette circonstance, et de l'existence même, comme de la silencieuse force, de cette identification à mon père, il y a quatre ans seulement (au cours de la méditation sur mon enfance et sur ma vie qui a suivi celle d'août 79 à mars 80 sur mes parents) Cette identification était comme le coeur paisible et puissant d'une identification à la famille que nous formions, mes parents, ma soeur (qui était mon \unsure{aînée|aîné} de quatre ans) et moi. Je vouais une admiration et un amour sans limites à mon père comme à ma mère. Leur personne était pour moi la mesure de toutes choses.

Cela ne signifie nullement que mon attitude à leur égard était celle \unsure{d'une} approbation automatique, d'une admiration béate. Je ne savais sans doute pas qu'ils étaient la mesure de toute chose pour moi, mais je savais fort bien qu'il étaient faillibles comme moi, et il n'y avait en moi aucune crainte qui m'aurait empêché de constater un désaccord et de le manifester clairement\unsure{..|.} Dans les conflits qui m'entouraient\unsure{;|,} je ne craignais pas de prendre partie à ma façon. Ce ne touchait en rien à une certaine foi, à une assurance qui formaient l'assise profonde, inébranlable de mon être - plutôt, cela découlait spontanément de \unsure{cette} foi, de cette assurance même.

Il arrivait que mon père, dans des accès de colère impuissante alors que ma soeur (sans en avoir l'air) prenait plaisir à le provoquer, la frappe avec brutalité - et à chaque fois j'en étais outragé, dans un élan de solidarité sans réserves avec ma soeur. C'étaient là je crois les seuls gros nuages qui \unsure{passaien|passaient} dans ma relation à\unsure{:}mon père (il n'y en avait pas avec ma mère). Ce n'est pas que j'approuvais les tours parfois pendables de ma soeur, pas plus je crois qu'ils ne me troublaient vraiment - ce n'était pas \mbox{e l l e} qui était pour moi la mesure des choses. Ses tours (dont la raison sûrement m'échappait tout autant qu'à mon père, qui "marchait" à tous les coups, ou à ma mère qui n'avait garde d'intervenir ni avant ni après) - ces tours en un sens ne tiraient pas vraiment à conséquence pour moi. C'était ma soeur, elle était comme elle était, voilà \unsure{tou|tout} Mais que \mbox{m o n}\quad\mbox{p è r e} se laisse aller à une telle brutalité aveugle...

Les trois êtres les plus proches, qui ensemble ont constitué comme la matrice de mes premières années, étaient déchirés par le conflit, opposant chacun

\noindent\rule{5em}{0.4pt}\par
(*) Voir la note "Le massacre", n° 87.

% archive sheet number: 130
